\documentclass{book}

\title{Parsing Foundations}
\newcommand{\booksubtitle}{A Study of Data and Parsing}
\newcommand{\booklicense}{Creative Commons Zero 1.0 Universal}

\author{Richard R. Masters}
\newcommand{\authorsubtitle}{Seattle, WA}

% Create convenient commands \booktitle and \bookauthor
\makeatletter
\newcommand{\booktitle}{\@title}
\newcommand{\bookauthor}{\@author}
\makeatother

% This utf8 declaration is not needed for versions of latex > 2018 but may
% be helpful for older software. Eventually it may not be worth keeping.
\usepackage[utf8]{inputenc}  
\usepackage{fix-cm}  % this package allows large \fontsize
\usepackage{tikz}    % this is for graphics. e.g. rectangle on title page
\usepackage{amsmath} % Used by equations

% The following dimensions specify 4.75" X 7.5" content on 6 3/8" by 9 1/4"
% paper. The paper width and height can be tweaked as required and the content
% should size to fit within the margins accordingly.
%
% The (inside) bindingoffset should be larger for books with more pages. Some
% standard recommended sizes are .375in minimum up to 1in for 600+ page books.
% Sizes .75in and .875in are also recommended roughly at 150 and 400 pages.
\usepackage[bindingoffset=0.625in,
            left=.5in, right=.5in,
            top=.8125in, bottom=.9375in,
            paperwidth=6.375in, paperheight=9.25in]{geometry}
% Here is an alternative geometry for reading on letter size paper:
% \usepackage[margin=.75in, paperwidth=8.5in, paperheight=11in]{geometry}

\renewcommand{\contentsname}{Table of Contents} % default is {Contents}
\usepackage{makeidx}
\makeindex % Initialize an index so we can add entries with \index

% Content Starts Here
\begin{document}
\frontmatter

% ---- Half Title Page ----
% current geometry will be restored after title page
\newgeometry{top=1.75in,bottom=.5in}
\begin{titlepage}
\begin{flushleft}

% Title
\textbf{\fontfamily{qcs}\fontsize{48}{54}\selectfont Parsing\\Foundations\\}

% Draw a line 4pt high
\par\noindent\rule{\textwidth}{4pt}\\

% Subtitle
% Shaded box from left to right. Text node is midway (centered).
\begin{tikzpicture}
\shade[bottom color=lightgray,top color=white]
    (0,0) rectangle (\textwidth, 1.5)
    node[midway] {\textbf{\large \textit{\booksubtitle}}};
\end{tikzpicture}

% Edition Number
\begin{flushright}
\Large Unfinished Draft of First Edition
\end{flushright}

% \vspace{\fill}
\vspace{\fill}

\end{flushleft}
\end{titlepage}
\restoregeometry
% ---- End of Half Title Page ----


% No page numbers on the Frontispiece page
\thispagestyle{empty}


% ---- Title Page ----
% current geometry will be restored after title page
\newgeometry{top=1.75in,bottom=.5in}
\begin{titlepage}
\begin{flushleft}

% Title
\textbf{\fontfamily{qcs}\fontsize{48}{54}\selectfont Parsing\\Foundations\\}

% Draw a line 4pt high
\par\noindent\rule{\textwidth}{4pt}\\

% Shaded box from left to right with Subtitle
% The text node is midway (centered).
\begin{tikzpicture}
\shade[bottom color=lightgray,top color=white]
    (0,0) rectangle (\textwidth, 1.5)
    node[midway] {\textbf{\large \textit{\booksubtitle}}};
\end{tikzpicture}

% Edition Number
\begin{flushright}
\Large Unfinished Draft of First Edition
\end{flushright}

\vspace{\fill}

% Author and Location
\textbf{\large \bookauthor}\\[3.5pt]
\textbf{\large \textit{\authorsubtitle}}

\vspace{\fill}

% Self Publishing Logo. Free to use: CC0 license.
% The source file is book.svg. If you change the svg, you must then convert
% it to pdf. There are many online and offline tools available to do that.
\begin{center}
\includegraphics{booksvg.pdf}\\[4pt]
\fontfamily{lmtt}\small{Self Publishers Worldwide\\
Seattle San Francisco New York\\
London Paris Rome Beijing Barcelona}
\end{center}

\end{flushleft}
\end{titlepage}
\restoregeometry
% ---- End of Title Page ----

% Do not show page numbers on colophon page
\thispagestyle{empty}

\begin{flushleft}
\vspace*{\fill}
This book was typeset using \LaTeX{} software.\\
\vspace{\fill}
Copyright \textcopyright{} \the\year{}  \bookauthor\\
License: \booklicense
\end{flushleft}

% A title page resets the page # to 1, but the second title page
% was actually page 3. So add two to page counter.
\addtocounter{page}{2}

% The asterisk excludes chapter from the table of contents.
\chapter*{Preface}
This book was created from a desire to make parsing easier to understand.
The author struggled to learn a variety of topics in the field parsing.
It was not that the ideas themselves were particularly difficult, but that
it required a huge investment in both time and money to acquire, understand,
and organize the relevant information. After all that effort, it was clear
that parsing could be explained a lot better, and also that there is a lot
of existing research, some of it several decades old, that could benefit
from a wider audience.

\section*{Prior Work}
Prior to the writing of this book, there were not that many resources for
learning about parsing that were both freely available and of high-quality.
There were no more than a handful of books about parsing on the market and
the most relevant to a general audience was a relatively expensive book about
compilers. There are also quite a few papers and articles written for academic
journals which may not be accessible or understandable to a general audience.
Finally, while there are some good online resources, they are lost within a sea
of information that is either outdated or poorly produced, or only covers a
limited set of topics. 

In contrast, the goal of this book is to be accessible and as easy to understand
as possible. This book is built from the ground up with relevant topics,
starting from the basics and building up from there, 
while using a simple and intuitive style of presentation.

Some of the explanations in this book may seem excessively long so it is
worth noting that keeping explanations short is a much lower priority
for this book than providing a thorough explanation. The idea is that far
more time is wasted by avoiding difficult material, skimming it repeatedly,
and finally focusing your mind enough to understand it, than is wasted
reading and understanding material that is excessively explained.


\section*{Audience}
This book can serve as an introduction to topics on data and parsing for
the working software developer or student of computer science. The book
covers foundational topics of parsing using concrete examples and common
tools used in the field. In an academic setting, the book would serve well as
a companion book for courses that only cover parsing as part of a larger
topic, such as compiling.

\section*{Uses for this Book}
The book may also be of interest to you if you are an educator, a blogger,
or if you are writing a book. This document is licensed under the Creative
Commons Zero v1.0 Universal license. You are free to copy this content,
distribute it, include it in your own content, or sell it, and you may do so
\textit{without attribution}. Just copy what you want and move on.  Read the
license text at https://creativecommons.org/publicdomain/zero/1.0/legalcode
for more details.

\section*{Prerequisites}
For the chapters on data, the reader should have a basic knowledge of
computers, such as knowing about a monitor, mouse, disk, file, and an
application.

% Three-level Table of Contents
\setcounter{tocdepth}{3}
\tableofcontents

\mainmatter

\chapter{Introduction}
This book is about data and parsing in the context of computing. Data is
defined by Wikipedia as any sequence of symbols for which "a specific act of
interpretation" gives those symbols meaning. Parsing is an act of
interpreting data based on how it matches structural patterns.

\backmatter
\addcontentsline{toc}{chapter}{Index}
\printindex
\end{document}
